\chapter*{怎样读这本书}
\addcontentsline{toc}{chapter}{怎样读这本书}
\markboth{怎样读这本书}{}

\section*{三层阅读标记}

同一章里，内容分成三层，用颜色标出来。颜色只是人为的路标：

\begin{itemize}[itemsep=4pt]
\item \textbf{故事主线（无标记的正文）}：只要求初中数学和常识。从头读到尾不跳步，就能理解核心因果。这是全书的骨架。
\item \textbf{【原理桥梁】绿色框}：用一点比例、图像、简单代数和概率，把定性的“是这样”变成定量的“是多少”。多花几分钟，理解会深一个档次。
\item \textbf{【深入挑战】紫色框}：通往更高层的窗口：轨道、自由能、反应机理、统计模型。第一次读可以整框跳过，绝不影响主线；哪天好奇心发作了再回来。
\end{itemize}

\section*{每章的走法}

每一章都从你身边拿起一样东西——一勺盐、一颗鸡蛋、一片叶子——然后沿着固定的路线往下挖：先看现象，再钻进粒子层，接着找规则，最后才轮到符号登场。章尾我们会回到开头那样东西，用新学到的模型重新解释它。如果你发现开头那个“理所当然”的现象忽然变得意味深长，这一章就成功了。

章与章之间还有反复登场的“老朋友”：一片面包、一颗苹果、一块电池、一块肥皂、一粒药片、一片叶子、一滴血、一个细菌。它们在元素章、化学键章、代谢章、基因章里一次次回来，每次都被看得更深一层。遇到老朋友时，不妨回想一下：上次见到它时，你理解到了哪一层？

\section*{五个贯穿问题}

面对任何物质或生命过程，我们都会依次追问五个问题：

\begin{enumerate}[itemsep=2pt]
\item 它由什么组成？
\item 这些组成部分怎样连接和排列？
\item 什么能量和概率规则决定它能否变化？
\item 变化有多快，又怎样受到控制？
\item 在生命中，这个过程怎样被保存、复制、调节和选择？
\end{enumerate}

读完全书后，希望你面对新闻里的一种新药、货架上的一种饮料、身体里的一次发烧，也能自己问出这五个问题。那时，这本书的方法就真正归你了。

\section*{动手与不动手}

书里有一些【动手试一试】栏目，全部只用厨房、药店和文具店能安全买到的东西。每个实验都写了安全提示，请照着做。还有一些现象——碱金属遇水、放射性衰变、细菌培养——我们只会描述或用视频资料示意，并明确告诉你：不要在家尝试。好奇心值得鼓励，但好的科学家首先懂得保护自己和别人。
